% ============================================================
% section4_processing.tex - Раздел 4: Оконная фильтрация и БПФ
% ============================================================
% !TeX program = xelatex
% Компиляция: xelatex --shell-escape main.tex

\section{Цифровая обработка: Окна и спектры}

В предыдущем разделе мы успешно загрузили и структурировали наши временные ТГц-импульсы. Однако фундаментальные физические свойства исследуемой проволочной решетки (такие как частотный коэффициент пропускания по мощности) проявляются в частотной (спектральной) области. 

Для перехода из временной области в частотную в цифровой обработке сигналов (DSP — Digital Signal Processing) используется классический инструмент — Быстрое преобразование Фурье (БПФ, или FFT — Fast Fourier Transform). Наша задача — использовать БПФ как готовый инженерный рецепт для превращения зависимости напряженности поля от времени в зависимость амплитуды от частоты.

\note{\textbf{Для любознательных: Математическая теория спектрального анализа}\\
В этой книге мы не будем углубляться в строгое математическое доказательство рядов и интегралов Фурье, свойства ортогональности гармонических функций и теоремы дискретизации. Мы рассматриваем БПФ исключительно как практический инструмент. Если ты хочешь глубоко изучить теоретические основы спектрального анализа и цифровой фильтрации, мы рекомендуем обратиться к классическим учебникам, например: \textbf{А. Оппенгейм, Р. Шафер. Цифровая обработка сигналов} или \textbf{С. Смит. Цифровая обработка сигналов. Практическое руководство для инженеров и научных работников}.}

Давай разберем пошаговый рецепт подготовки сигнала и расчета спектра, а затем оформим его в виде готового модуля \texttt{spectrum.py}.

\subsection{Формат данных и типы в NumPy: Знакомимся с np.ndarray}

Прежде чем приступать к цифровой обработке, давай вспомним, в каком именно формате возвращает данные наша функция \texttt{load\_tds\_data}, написанная в предыдущем разделе. При вызове:
\begin{minted}{python}
t, E = load_tds_data(filepath)
\end{minted}
мы получаем два связанных одномерных числовых вектора:
\begin{enumerate}
    \item \texttt{t} (массив временных отсчетов) — представляет собой сетку времени в пикосекундах ($\text{пс}$), обычно с шагом порядка $0.05$--$0.1$~пс, общей длительностью записи около $100$~пс.
    \item \texttt{E} (массив амплитуды поля) — содержит соответствующие значения напряженности электрического поля терагерцового импульса в условных единицах (у.е.).
\end{enumerate}

Оба этих объекта имеют тип \texttt{np.ndarray} (аббревиатура от \textit{N-dimensional array} — многомерный массив NumPy). Это фундаментальный тип данных в научных вычислениях на Python. В отличие от стандартных списков Python (\texttt{list}), которые могут содержать объекты любых типов и хранятся в памяти в виде разбросанных ссылок, массив \texttt{np.ndarray} устроен принципиально иначе:
\begin{itemize}
    \item \textbf{Однородность (homogeneous):} все элементы массива строго одного типа (обычно вещественные числа двойной точности \texttt{float64}).
    \item \textbf{Непрерывность в памяти (C-contiguous):} элементы лежат в оперативной памяти единым непрерывным блоком, как в языках C или Fortran.
    \item \textbf{Векторизация (vectorization):} математические операции над массивами выполняются на уровне скомпилированного C-кода без медленных циклов Python. Например, выражение \texttt{E - 5.0} мгновенно вычтет число $5.0$ из всех элементов массива параллельно.
\end{itemize}

Именно благодаря свойствам \texttt{np.ndarray} мы можем осуществлять сложнейшую обработку ТГц-сигналов (вычисление средних, наложение окон, преобразование Фурье) в реальном времени со скоростью сотен кадров в секунду!

\subsection{Шаг 1: Удаление постоянного смещения (DC offset)}

Любой реальный прибор вносит в измеряемый сигнал небольшое аппаратное смещение нуля (постоянную составляющую, или DC offset) — постоянный вертикальный сдвиг всей ТГц-осциллограммы вверх или вниз относительно истинного нуля. Этот сдвиг возникает из-за тепловых наводок в приемной электронике.

Если применить БПФ к сигналу с таким смещением, алгоритм воспримет этот сдвиг как волну нулевой частоты. В результате на спектре возникнет ложный гигантский пик на частоте $0$~ТГц, который скроет от нас полезные данные на низких частотах.

Рецепт исправления предельно прост: нужно вычислить среднее значение сигнала и вычесть его из каждой точки:

\begin{minipage}{\linewidth}
	\captionof{listing}{Удаление постоянной составляющей (spectrum.py)}
	\begin{minted}{python}
import numpy as np

def remove_dc(E: np.ndarray) -> np.ndarray:
    """Вычитает среднее арифметическое значение из массива амплитуд поля."""
    return E - np.mean(E)
	\end{minted}
\end{minipage}

Обрати внимание на синтаксис определения функции \texttt{remove\_dc(E: np.ndarray) -> np.ndarray}. Здесь впервые в нашей книге используются \textbf{аннотации типов} (type hinting / тайп-хинтинг — подсказки типов), ставшие стандартом в современном Python:
\begin{itemize}
    \item Двоеточие после имени аргумента (\texttt{E: np.ndarray}) указывает, что функция ожидает на вход массив NumPy.
    \item Стрелка (\texttt{-> np.ndarray}) перед двоеточием в конце строки декларирует, что функция возвращает тоже массив NumPy.
\end{itemize}
Сами по себе аннотации типов не замедляют выполнение программы и не вызывают ошибок при запуске, если передать не тот тип (Python остается языком с динамической типизацией). Однако они критически важны для разработчика: современные редакторы кода (IDE) сразу же подсветят ошибку в коде, если ты попытаешься передать в функцию обычную строку или список вместо массива, а также предложат автозаполнение доступных методов объекта.

\subsection{Шаг 2: Оконная фильтрация (Гауссово окно)}

Терагерцовый импульс представляет собой очень короткий всплеск поля длительностью менее $1$--$2$~пс. Однако запись сигнала длится значительно дольше (например, порядка $100$~пс). В «хвосте» сигнала полезного импульса уже нет — там регистрируются только случайные шумы прибора и эхо-сигналы (паразитные переотражения ТГц-волны внутри деталей прибора).

Чтобы аккуратно вырезать полезный импульс и плавно подавить шум на краях, мы применим оконную фильтрацию — умножим наш временной сигнал на оконную функцию $W(t)$ (Гауссово окно), которая описывается классическим уравнением:
\begin{equation}
W(t) = \exp\left( -\frac{(t - t_{\text{peak}})^2}{2\sigma^2} \right)
\end{equation}
где $t_{\text{peak}}$ — положение временного центра (пика) ТГц-импульса, а $\sigma$ — полуширина окна, определяющая скорость затухания функции к краям записи. 

Ширину этого окна во временной области мы зададим в файле конфигурации \texttt{config.py} параметром \texttt{SIGMA\_PS}. По умолчанию мы установим стандартное значение $\sigma = 20.0$~пс, что позволит аккуратно локализовать ТГц-импульс.

Перед генерацией окна необходимо автоматически найти положение главного пика импульса $t_{\text{peak}}$, чтобы точно отцентрировать окно по временной шкале:

\begin{minipage}{\linewidth}
	\captionof{listing}{Центрирование и наложение Гауссова окна (spectrum.py)}
	\begin{minted}{python}
import config

def find_peak_center(E: np.ndarray) -> int:
    """Возвращает индекс максимального по модулю значения (центр импульса)."""
    return int(np.argmax(np.abs(E)))

def apply_gaussian_window(t: np.ndarray, E: np.ndarray, sigma_ps: float = None) -> tuple[np.ndarray, np.ndarray]:
    """
    Создает Гауссово окно с заданной шириной sigma_ps (в пикосекундах) 
    вокруг главного пика и плавно сводит сигнал к нулю по краям.
    """
    if sigma_ps is None:
        sigma_ps = config.SIGMA_PS
        
    peak_idx = find_peak_center(E)
    t_peak = t[peak_idx]
    
    # Математическая формула Гауссовой кривой: exp(- (t - t_peak)^2 / (2 * sigma^2))
    window = np.exp(-0.5 * ((t - t_peak) / sigma_ps) ** 2)
    
    E_windowed = E * window
    return E_windowed, window
	\end{minted}
\end{minipage}

Давай внимательно разберем заголовок и структуру функции \texttt{apply\_gaussian\_window}. Здесь мы видим несколько новых для нас конструкций современного Python:
\begin{enumerate}
    \item \textbf{Необязательный аргумент со значением по умолчанию (\texttt{sigma\_ps: float = None}):} 
    Конструкция \texttt{= None} указывает, что параметр \texttt{sigma\_ps} является опциональным (необязательным для передачи). Если при вызове функции мы не укажем ширину окна (например, вызовем \texttt{apply\_gaussian\_window(t, E)}), Python автоматически присвоит переменной \texttt{sigma\_ps} специальное пустое значение \texttt{None}. Внутри функции мы проверяем это условие:
    \begin{minted}{python}
    if sigma_ps is None:
        sigma_ps = config.SIGMA_PS
    \end{minted}
    и подставляем дефолтное (заданное по умолчанию) значение из нашего конфигурационного файла. Это очень красивый паттерн проектирования: функция остается гибкой (мы можем передать любую ширину), но при этом имеет надежное стандартное поведение «из коробки».
    
    \item \textbf{Аннотация возвращаемого типа \texttt{-> tuple[np.ndarray, np.ndarray]}:}
    Слово \texttt{tuple} (кортеж / упорядоченный неизменяемый набор данных) указывает на то, что функция возвращает не одиночное значение, а кортеж из двух элементов. В квадратных скобках \texttt{[np.ndarray, np.ndarray]} подробно расписано, что и первый, и второй элементы этого кортежа являются массивами NumPy. В конце функции мы видим строчку:
    \begin{minted}{python}
    return E_windowed, window
    \end{minted}
    Это позволяет нам при вызове функции элегантно распаковать результат в две независимые переменные одной строкой:
    \begin{minted}{python}
    E_filtered, W = apply_gaussian_window(t, E)
    \end{minted}
\end{enumerate}

Кроме борьбы с эхо-сигналами, у оконного фильтра есть еще один важный практический эффект — контролируемое сглаживание спектра. В реальном ТГц-эксперименте на спектрах часто видны очень резкие и глубокие провалы — узкие линии поглощения (например, молекулами водяного пара в воздухе). Применение Гауссова окна сглаживает спектральную кривую: узкие линии поглощения и мелкие шумовые осцилляции становятся менее глубокими и размываются. Это помогает убрать лишнюю «изрезанность» графика и получить более чистую, плавную огибающую спектра пропускания нашей решетки.

\warning{Вычитание постоянной составляющей (DC offset) обязательно должно выполняться \textbf{до} наложения окна! Если сначала умножить сигнал со смещением на окно, постоянный сдвиг превратится в плавный Гауссов «горб» во временной области. Такой искусственный сигнал уже невозможно отфильтровать простым вычитанием константы.}

\subsection{Шаг 3: Расчет быстрого преобразования Фурье (БПФ)}

Теперь мы готовы перевести временной ТГц-импульс в спектр. В Python Быстрое преобразование Фурье реализовано в стандартной библиотеке \texttt{scipy.fft}. Так как наше поле $E(t)$ состоит из чисто вещественных (действительных) чисел, мы используем специализированную функцию \texttt{rfft} (Real FFT) вместо стандартной комплексной \texttt{fft}. Это экономит вычислительные ресурсы и автоматически убирает избыточные отрицательные частоты.

\note{\textbf{Важная деталь о названии «Real FFT»:} у начинающих программистов часто возникает путаница. Префикс \texttt{r} (от \textit{Real} — вещественный) в названии функции \texttt{rfft} относится исключительно к \textbf{входным данным} (то есть на вход подается массив вещественных чисел \texttt{E}, у которых мнимая часть равна нулю). Однако сам \textbf{результат работы функции} (выходной спектр) остается \textbf{комплексным}! Каждая спектральная гармоника по-прежнему имеет как амплитуду, так и фазу, поэтому возвращаемый массив \texttt{spectrum\_complex} состоит из комплексных чисел (\texttt{complex128}). Чтобы избавиться от фазовой составляющей и получить привычный вещественный амплитудный спектр, мы обязаны взять модуль комплексных чисел с помощью функции \texttt{np.abs(spectrum\_complex)}.}

Частотную шкалу мы получим с помощью функции \texttt{rfftfreq}:

\begin{minipage}{\linewidth}
	\captionof{listing}{Вычисление спектра амплитуд (spectrum.py)}
	\begin{minted}{python}
from scipy.fft import rfft, rfftfreq

def compute_spectrum(t: np.ndarray, E: np.ndarray) -> tuple[np.ndarray, np.ndarray]:
    """
    Вычисляет спектр амплитуд ТГц-импульса.
    Возвращает ось частот (в ТГц) и вещественный амплитудный спектр.
    """
    N = len(E)
    dt = t[1] - t[0]  # Шаг временной сетки в пикосекундах
    
    # Вычисляем преобразование Фурье для вещественного сигнала
    spectrum_complex = rfft(E)
    
    # Модуль комплексного числа дает амплитуду спектральной гармоники
    amplitude_spectrum = np.abs(spectrum_complex)
    
    # Генерация шкалы частот.
    # Так как dt задан в пикосекундах, частоты freqs будут в ТГц
    freqs = rfftfreq(N, d=dt)
    
    return freqs, amplitude_spectrum
	\end{minted}
\end{minipage}

\subsection{Шаг 4: Усреднение сигналов в спектральной области}

При проведении ТГц-эксперимента для каждого угла поляризатора обычно записывают несколько повторных измерений (проходов) сигнала, чтобы повысить точность.

\warning{Ни в коем случае не усредняй исходные временные сигналы нескольких проходов перед расчетом БПФ! \\
\\
Из-за незначительного случайного временного дрожания прибора (фазового джиттера) ТГц-импульсы в разных проходах могут быть микроскопически сдвинуты друг относительно друга по оси времени. Если их усреднить во временной области, они частично погасят друг друга. На итоговом спектре это приведет к ложному падению амплитуд на высоких частотах.\\
\\
\textbf{Правильный рецепт:} сначала рассчитай БПФ-спектр для каждого прохода по отдельности, возьми модуль (амплитуду), и только после этого усредняй полученные амплитудные графики на частотной сетке.}

\subsection{Шаг 5: Вычисление пропускания по мощности}

Коэффициент пропускания $T(\nu)$ исследуемой решетки по мощности (power transmission) в ТГц-диапазоне показывает, какая часть мощности излучения прошла сквозь нее на каждой частоте. Мощность электромагнитной волны пропорциональна квадрату ее амплитуды. Поэтому формула пропускания выглядит следующим образом:
\begin{equation}
T(\nu) = \left| \frac{E_{\text{sample}}(\nu)}{E_{\text{bg}}(\nu)} \right|^2 = \left( \frac{A_{\text{sample}}(\nu)}{A_{\text{bg}}(\nu)} \right)^2
\end{equation}

Напишем функцию \texttt{calculate\_transmission}, которая объединяет все шаги в единый конвейер:

\begin{minipage}{\linewidth}
	\captionof{listing}{Конвейер расчета спектра пропускания (spectrum.py)}
	\begin{minted}{python}
def calculate_transmission(t_sig: np.ndarray, E_sig: np.ndarray, 
                           t_bg: np.ndarray, E_bg: np.ndarray) -> tuple[np.ndarray, np.ndarray, np.ndarray, np.ndarray]:
    """Полный цикл обработки: от сырого ТГц-сигнала до спектра пропускания по мощности."""
    
    # 1. Удаление постоянного смещения
    E_sig_clean = remove_dc(E_sig)
    E_bg_clean = remove_dc(E_bg)
    
    # 2. Наложение временного Гауссова окна
    E_sig_win, _ = apply_gaussian_window(t_sig, E_sig_clean)
    E_bg_win, _ = apply_gaussian_window(t_bg, E_bg_clean)
    
    # 3. Расчет БПФ
    freqs, spec_sig = compute_spectrum(t_sig, E_sig_win)
    _, spec_bg = compute_spectrum(t_bg, E_bg_win)
    
    # 4. Расчет пропускания по мощности (T = |E_sig / E_bg|^2)
    # Используем np.maximum для предотвращения деления на ноль на краях диапазона
    spec_bg_safe = np.maximum(spec_bg, 1e-10)
    transmission = (spec_sig / spec_bg_safe) ** 2
    
    return freqs, spec_sig, spec_bg, transmission
	\end{minted}
\end{minipage}

\begin{minitaskbox}
	\textbf{Мини-задание 2: Отображение динамического спектра в GUI (файл \texttt{main.py})} \\
	Для того чтобы наглядно увидеть, как оконная фильтрация влияет на спектр, тебе предстоит модифицировать Шаг 1 графического интерфейса в файле \texttt{main.py}. На этом шаге пользователь настраивает параметры Гауссова окна.
	\\ \\
	\textbf{Твоя задача:} добавить на окно настроек под графиком временного сигнала вторую координатную ось \texttt{ax2} (связанную по горизонтали). При изменении положения интерактивного ползунка ширины окна $\sigma$ программа должна автоматически:
	\begin{enumerate}
		\item Пересчитывать и строить заокненный временной ТГц-сигнал на верхней панели.
		\item Вычислять БПФ-спектр этого сигнала с помощью \texttt{compute\_spectrum} и в реальном времени строить его на нижней панели (ось Y должна быть логарифмической в дБ: \texttt{20 * np.log10(spec)}).
	\end{enumerate}
	Это наглядно покажет, как сужение временного окна $\sigma_t$ сглаживает спектр и подавляет паразитные частотные осцилляции от эхо-сигналов.
\end{minitaskbox}

\subsection{Полный листинг модуля}

Полный, готовый к работе исходный код расчетного модуля \texttt{spectrum.py} (с решенными проверочными блоками и автономными тестами на синтетических сигналах) вынесен в \textbf{\textit{Приложение~\ref{app:spectrum}}}. Ты всегда можешь использовать его для сверки со своим кодом.

\subsection{Итоги раздела}
В этом разделе мы научились применять базовые цифровые методы для обработки экспериментальных ТГц-сигналов:
\begin{enumerate}
	\item Очистили сигнал от постоянного вертикального смещения прибора с помощью функции \texttt{remove\_dc}.
	\item Научились плавно вырезать полезную часть импульса и подавлять шумы на краях записи с помощью Гауссова окна \texttt{apply\_gaussian\_window}.
	\item Использовали инструмент быстрого преобразования Фурье \texttt{compute\_spectrum} для мгновенного перехода в частотную область.
	\item Объединили все операции в конвейер \texttt{calculate\_transmission} для расчета экспериментального спектра пропускания по мощности.
\end{enumerate}

Теперь, когда у нас готов расчетный фундамент, в следующем разделе мы применим его ко всему массиву экспериментальных данных и построим красивое угловое семейство спектров пропускания!
